\section{Methods}
\label{sec:methods}

\subsection{Heuristic Categorization Scheme}
We utilize a heuristic categorization scheme to structure our analysis of natural molecular hallucinations (C3). Each type is detectable through rule-based classification using RDKit~\cite{landrum2023rdkit} as the ground-truth engine. We emphasize that this scheme is a heuristic tool, not an exhaustive semantic ontology. The keyword sets were intentionally constrained to maintain interpretability and reproducibility rather than maximize coverage. The categories are mutually exclusive and applied in priority order (SP $>$ PI $>$ MF $>$ CC):

\begin{itemize}
    \item \textbf{Structural Phantom (SP):} The generated text references chemical elements or functional groups that are absent from the molecular structure. Detection: for a predefined heuristic dictionary of element keywords (e.g., ``fluorine'', ``fluoro''), we check whether the corresponding element is present in the molecule's atom set via \texttt{mol.GetAtoms()}. If a keyword matches but the element is absent, the description is classified as SP.
    \item \textbf{Property Inversion (PI):} The generated text claims physicochemical properties that contradict computed descriptors. Detection: we compute MolLogP via RDKit. If LogP $> 3$ and the text contains keywords from a predefined hydrophilic dictionary, or LogP $< 1$ and the text contains keywords from a lipophilic dictionary, the description is classified as PI.
    \item \textbf{Mechanism Fabrication (MF):} The generated text names specific biological targets or pathways. Detection: exact keyword matching against a fixed list of 9 common targets (e.g., ACE2, HER2, EGFR).
    \item \textbf{Contextual Confabulation (CC):} The residual category. Descriptions that do not trigger SP, PI, or MF criteria are classified as CC. This category represents scientifically plausible but generic narrative hallucination.
\end{itemize}

We acknowledge that CC functions as a residual class and encompasses heterogeneous subtypes that future work should further differentiate.

\subsection{Experimental Conditions}
Our ablation framework consists of nine conditions, designed to isolate the contributions of semantic content, scientific register, and hallucination type:

\begin{itemize}
    \item \textbf{C0 (Baseline):} SMILES notation only, with no augmentation.
    \item \textbf{C1 (Factual):} SMILES augmented with a factual description generated by RDKit (molecular weight, LogP, hydrogen bond donors/acceptors, TPSA, ring count, atom composition).
    \item \textbf{C2 (Chemical Priming):} SMILES augmented with a template-generated text using scientific-sounding but semantically vacuous language that includes chemical vocabulary. We note that C2 is \emph{not} a pure stylistic control because it contains domain-relevant vocabulary; C2b serves as the stricter control.
    \item \textbf{C2b (Pure Gibberish):} SMILES augmented with template-generated text using exclusively non-chemical vocabulary. This condition strictly controls for prompt length and linguistic complexity without domain priming.
    \item \textbf{C3 (Free Hallucination):} SMILES augmented with an unconstrained LLM-generated description, replicating the protocol of prior work~\cite{chen2023llm4mol}.
    \item \textbf{C4a (Structural Augmentation):} SMILES augmented with text generated via prompt-constrained forcing to explicitly describe structural features and functional groups. We emphasize that C4a--c represent prompt-constrained semantic perturbations rather than naturally occurring hallucinations classified post-hoc. To validate this constraint, we applied the heuristic classifier to a representative sample of generated C4a descriptions ($n=49$). The results revealed a notable methodological limitation: 0\% of C4a outputs were classified as SP, and the vast majority defaulted to CC. This occurred because the LLM generated scientifically plausible structural descriptions (e.g., ``contains a benzene ring'') that do not reference absent elements---the specific criterion required by the SP heuristic (which checks for keywords like ``fluorine'' or ``chloro'' against the molecule's actual atom set). Consequently, the C4a--c conditions must be strictly interpreted as \textit{prompt-constrained semantic perturbations} (focusing the LLM on structural, property, or mechanism topics), rather than verified taxonomy classes.
    \item \textbf{C4b (Property Inversion):} SMILES augmented with text generated via prompt-constrained forcing to explicitly describe solubility, lipophilicity, and permeability.
    \item \textbf{C4c (Mechanism Fabrication):} SMILES augmented with text generated via prompt-constrained forcing to explicitly hypothesize biological targets and cellular mechanisms.
    \item \textbf{C5 (Random-Permutation):} SMILES augmented with a hallucination originally generated for a \emph{different} molecule, assigned via a fixed random permutation ($seed=42$). This preserves the scientific register and statistical properties of hallucinated text while completely disrupting semantic alignment with the target molecule.
\end{itemize}

\subsection{Prediction Protocol}
All prediction queries use a fixed zero-shot prompt template:

\begin{quote}
\texttt{Given this molecule (SMILES: \{smiles\})}\\
\texttt{[Additional information: \{augmentation\}]}\\
\texttt{On a scale of 0 to 1, what is the probability that this molecule can \{task\}?}\\
\texttt{Respond with ONLY a number between 0 and 1. Nothing else.}
\end{quote}

The augmentation line is omitted for C0. The task description is ``penetrate the blood-brain barrier'' for BBBP and ``inhibit the BACE-1 enzyme'' for BACE. For the extended generalization benchmarks, the task descriptions were: ``exhibit toxicity in the NR-AR assay'' for Tox21 (selected as a representative nuclear receptor assay to provide a mechanistic contrast to the enzymatic reasoning required for BACE), and ``provide a water solubility (logS) value consistent with the molecule's structure'' for ESOL. For ESOL, the task was framed as log-solubility regression; RMSE was calculated directly on the predicted numerical values extracted from LLM responses. All other prompt elements remained identical to the primary ablation. The numerical output was deterministically extracted, with malformed responses defaulting to 0.5 for classification tasks and 0.0 for regression.

\subsection{Implementation Details}
\label{sec:implementation}
\textbf{Model.} Primary ablation experiments were conducted with \texttt{Llama-\allowbreak 3.1-\allowbreak 8B-\allowbreak Instant} via the Groq API. Cross-model validation was extended to \texttt{Llama-\allowbreak 3.3-\allowbreak 70B-\allowbreak Versatile}, as described in Section~\ref{sec:implementation_robustness}.

\textbf{Generation parameters.} Prediction queries: temperature $= 0.0$, $\texttt{max\_tokens} = 10$, $\texttt{seed} = 42$. Hallucination generation: temperature $= 0.9$, $\texttt{max\_tokens} = 150$, no fixed seed. We explicitly distinguish between stochastic hallucination generation (no fixed seed) and quasi-deterministic prediction inference (seed=42), noting that API-level non-determinism may still introduce minor variability. To ensure exact reproducibility despite stochastic generation, the static dataset containing all generated augmentations used for evaluation is provided in the repository. Failed API calls (rate limits, parsing errors) defaulted to $p = 0.5$; the proportion of such defaults is reported in Appendix~\ref{subsec:api_failures}.

\textbf{Data.} Datasets were loaded via DeepChem~\cite{ramsundar2019deepchem}. Both BBBP ($N=204$) and BACE ($N=152$) were partitioned using scaffold splitting~\cite{bemis1996properties} via DeepChem's built-in splitter (80/10/10 train/valid/test). Only the held-out test set was used for evaluation.

\textbf{Factual descriptions.} For condition C1, factual descriptions were generated using RDKit molecular descriptors: molecular weight, Wildman-Crippen LogP, number of hydrogen bond donors and acceptors, topological polar surface area, heavy atom count, and ring count.

\textbf{Traditional ML Baseline.} To provide context for the zero-shot LLM predictions, we trained a Random Forest classifier (500 trees) using Morgan fingerprints (radius=2, 2048 bits) on the training and validation partitions of the scaffold split. The baseline was evaluated on the exact same test sets ($N=204$ for BBBP, $N=152$ for BACE) to ensure strict comparability.

\textbf{Extended Generalization Benchmarks.} To assess 
cross-task generalizability of the primary findings, a 
targeted subset of conditions---baseline (C0), free 
hallucination (C3), and prompt-constrained structural 
augmentation (C4a)---was evaluated on two additional 
MoleculeNet benchmarks: Tox21 (multi-task toxicity 
classification, evaluated on the first task column) and 
ESOL (aqueous solubility regression). These three 
conditions were selected to represent the no-augmentation 
reference, the maximal perturbation condition, and the 
best-performing augmentation identified in the primary 
ablation, respectively. Full nine-condition ablation on 
these datasets is reserved for future work given the 
exploratory scope of this study. Results are reported in 
Table~\ref{tab:generalization}.

\subsection{Evaluation and Calibration Metrics}
Discriminative performance was assessed using the area under the receiver operating characteristic curve (ROC-AUC). To account for the single-split evaluation, 95\% confidence intervals were computed via bootstrap resampling (1,000 iterations, seed $= 42$). 

To characterize the observed ``distributional compression'' and its impact on model reliability, we evaluate model calibration using the Brier score and Expected Calibration Error (ECE). The Brier score measures the mean squared difference between predicted probabilities ($f_i$) and actual outcomes ($y_i$):
\begin{equation}
BS = \frac{1}{N} \sum_{i=1}^{N} (f_i - y_i)^2
\end{equation}
ECE partitions the predicted probabilities into $M=10$ equally spaced bins ($B_m$) and computes the weighted average of the absolute difference between bin accuracy and bin confidence:
\begin{equation}
ECE = \sum_{m=1}^{M} \frac{|B_m|}{N} |acc(B_m) - conf(B_m)|
\end{equation}

Statistical significance for ROC-AUC differences relative to the baseline (C0) was assessed using a bootstrap paired difference test. Multiple comparisons were corrected using the Benjamini-Hochberg procedure (hereafter $p_{\text{BH}}$). We also report Cohen's $d$ on the raw prediction score distributions to quantify the shift in model confidence induced by semantic perturbations.

\subsection{Preliminary Scaling Exploration}
\label{sec:implementation_robustness}
To assess whether the observed phenomena persist across model scales, we performed two preliminary exploration steps:
\begin{enumerate}
    \item \textbf{Multi-Seed Stability (8B):} For the primary Llama-3.1-8B model, we repeated the prediction protocol for conditions C0 and C3 across three independent runs (seeds 42, 123, 777) to observe performance variance.
    \item \textbf{Pilot Scale Exploration (70B):} To observe behavior at a larger scale, we executed the prediction protocol for C0 and C3 on the \texttt{Llama-\allowbreak 3.3-\allowbreak 70B-\allowbreak Versatile} model.
\end{enumerate}
This exploration was designed to distinguish systematic semantic effects from model-specific quirks at the 8B scale.
